(1) 均匀极化给 \(\rho_b=-\nabla\cdot\boldsymbol P=0\)。球面外法向为 \(\hat r\)，故 \(\sigma_b(\theta)=\boldsymbol P\cdot\hat r=P_0\cos\theta\)。正负面电荷分布于相反半球。

(2) 极化分布产生的总偶极矩是 \(\boldsymbol p=\int_{\rm sphere}\boldsymbol P\,dV=(4\pi a^3/3)P_0\hat z\)。亦可直接积分表面面电荷的 \(\int\boldsymbol r\sigma_b\,dA\)，两法结果一致。

(3) 设内部电势为 \(\phi_{\rm in}=Ar\cos\theta\)，外部为 \(\phi_{\rm out}=B\cos\theta/r^2\)。球面电势连续给 \(B=Aa^3\)。法向电位移连续条件用 \(D_{\rm in}=\epsilon_0E_{\rm in}+P\)，\(D_{\rm out}=\epsilon_0E_{\rm out}\)：
\[
\epsilon_0(-A)+P_0=\epsilon_0(2B/a^3)=2\epsilon_0A.
\]
所以 \(A=P_0/(3\epsilon_0)\)，内部场
\[
\boldsymbol E_{\rm in}=-\nabla\phi_{\rm in}
=-\frac{\boldsymbol P}{3\epsilon_0},
\]
是均匀的去极化场。

(4) 外部势为 \(\phi_{\rm out}=P_0a^3\cos\theta/(3\epsilon_0r^2)=p\cos\theta/(4\pi\epsilon_0r^2)\)，恰是点偶极子势。总束缚电荷 \(Q_b=\int\sigma_b\,dA=P_0a^2\int\cos\theta\,d\Omega=0\)。
